\section{Greek version [Apostolos Syropoulos \& Anastasios Dimou]}

\subsection{Εισαγωγή} 
Ο Alain Matthes μας έχει συνηθίσει σε ενδιαφέροντα πακέτα για το \LaTeX\ , που είναι μάλιστα πολύ σχετικά με τα δικά μας προγράμματα, το στυλ και το ύφος τους. Ένα τέτοιο παράδειγμα είναι και το \texttt{tkz-tab}, που παρουσιάστηκε πέρυσι στο \texttt{https://tassosdimou.gr/variation-table}.

Το πακέτο \textsf{alterqcm} είναι ακόμη ένα πακέτο του Alain Matthes για το \LaTeX\, που θα μας βοηθήσει στη κατασκευή καλαίσθητων διαγωνισμάτων με ερωτήσεις πολλαπλής επιλογής και σωστού-λάθους.

Το \textsf{alterqcm}  τροποποιήθηκε από τους Απόστολο Συρόπουλο και Τάσσο Δήμου έτσι, ώστε να προσαρμοστεί στα δεδομένα του ελληνικού εκπαιδευτικού συστήματος. 

 Το άρθρο αναπτύσσει με λεπτομέρειες και πολλά παραδείγματα τις δυνατότητες του \textsf{alterqcm}. Δίνει οδηγίες για τη χρήση του και στο τέλος  θα δοθούν μερικά παραδείγματα διαγωνισμάτων.

\subsection{Εγκατάσταση του πακέτου}	
Θα υποδείξουμε έναν απλό τρόπο εγκατάστασης του πακέτου. Δημιουργούμε ένα φάκελο, στον οποίο θα αποθηκευτούν όλα τα αρχεία, που θα επεξεργαστούμε, μελετώντας το \textsf{alterqcm}. Με άλλα λόγια, στον φάκελο αυτόν αποθηκεύουμε τα αρχεία \texttt{.tex}, τις εικόνες που θα χρησιμοποιηθούν και το αρχείο \texttt{alterqcm.sty}, που θα κατεβάσουμε από τη διεύθυνση \texttt{https://ctan.org/pkg/alterqcm?lang=en}. Το πακέτο θα φορτωθεί με την επιλογή \texttt{greek}, δηλαδή θα δώσουμε την εντολή:
\begin{verbatim}
\usepackage[greek]{alterqcm}
\end{verbatim}
Όλα τα αρχεία θα έχουν την κλασσική δομή των αρχείων \texttt{.tex}.

Στο πρώτο μέρος, το προοίμιο, θα τοποθετήσουμε τα: 
\begin{verbatim}
\documentclass[11pt,a4paper]{article}
\usepackage{xltxtra}
\usepackage{xgreek}
\usepackage{mathtools}
\usepackage{amsthm}
\usepackage{amssymb}
\usepackage{unicode-math}
\usepackage{xkeyval}
\usepackage{multirow,longtable}
\usepackage[greek]{alterqcm}
\usepackage{tkz-tab}
%%%%%%%%%%%%%%%%%%%%%%%%%%%%%%%%%%%%%%%%%%%%%%%%%%%%%%%%%
\parindent=0pt
\setmainfont[Mapping=tex-text,Ligatures=Common]{Minion Pro}
\setmathfont[Scale=MatchUppercase]{Asana Math}
\end{verbatim}	

Apostolos Syropoulos, and Anastasios Dimou

\begin{alterqcm}
  \AQquestion{Ερώτηση}{%
  {Επιλογή 1},
  {Επιλογή 2},
  {Επιλογή 3}
  }
 \end{alterqcm}
 
 \begin{alterqcm}[language=greek,VF,lq=60mm]
 \AQquestion[]{Ισχύει ότι $(α+β)^2=α^2+β^2$}
 \AQquestion[]{Αν $α\cdot β\geq 0$, τότε $\sqrt{α\cdot β}=\sqrt{α}\cdot\sqrt{β}$ } 
 \AQquestion[]{Είναι $|α|=α,\,\text{για κάθε}\ 
  x\in\mathbb{R}$}
 \end{alterqcm}